% Multi-horizon macros, minted by experiments/multihorizon_tex.py from
% results/multihorizon/*.csv. Guarded so a second \input is a no-op.
\ifdefined\mhNDays\else% GENERATED FILE -- do not edit by hand.
% Written by experiments/multihorizon_tex.py from results/multihorizon/*.csv.
\newcommand{\mhAzeroMOneDMVsNaive}{$-$19.5}
\newcommand{\mhAzeroMOneQLIKE}{0.1837}
\newcommand{\mhAzeroMThreeDMVsNaive}{$-$19.6}
\newcommand{\mhAzeroMThreeQLIKE}{0.1761}
\newcommand{\mhAzeroMTwoDMVsNaive}{$-$19.4}
\newcommand{\mhAzeroMTwoQLIKE}{0.1814}
\newcommand{\mhAzeroMTwopooledDMVsNaive}{$-$19.4}
\newcommand{\mhAzeroMTwopooledQLIKE}{0.1814}
\newcommand{\mhAzeroSlopeOnestepEarly}{0.80}
\newcommand{\mhAzeroSlopeOnestepLate}{0.62}
\newcommand{\mhAzeroSlopeSofarEarly}{0.16}
\newcommand{\mhAzeroSlopeSofarLate}{0.43}
\newcommand{\mhBlkMOneDMVsNaive}{$-$19.9}
\newcommand{\mhBlkMOneQLIKE}{0.1740}
\newcommand{\mhBlkMThreeDMVsNaive}{$-$20.0}
\newcommand{\mhBlkMThreeQLIKE}{0.1674}
\newcommand{\mhBlkMTwoDMVsNaive}{$-$19.8}
\newcommand{\mhBlkMTwoQLIKE}{0.1719}
\newcommand{\mhBlkMTwopooledDMVsNaive}{$-$19.8}
\newcommand{\mhBlkMTwopooledQLIKE}{0.1719}
\newcommand{\mhBlkSlopeOnestepEarly}{0.78}
\newcommand{\mhBlkSlopeOnestepLate}{0.65}
\newcommand{\mhBlkSlopeSofarEarly}{0.17}
\newcommand{\mhBlkSlopeSofarLate}{0.39}
\newcommand{\mhDMBlkVsAzeroMax}{$-$2.8}
\newcommand{\mhDMBlkVsAzeroMin}{$-$6.3}
\newcommand{\mhHOneFiveZeroZeroAzeroMTwoQLIKE}{0.1327}
\newcommand{\mhHOneFiveZeroZeroBlkMTwoQLIKE}{0.1238}
\newcommand{\mhHOneFiveZeroZeroDMBlkVsAzero}{$-$6.3}
\newcommand{\mhHOneFourThreeZeroAzeroMTwoQLIKE}{0.1492}
\newcommand{\mhHOneFourThreeZeroBlkMTwoQLIKE}{0.1378}
\newcommand{\mhHOneFourThreeZeroDMBlkVsAzero}{$-$6.1}
\newcommand{\mhHOneFourZeroZeroAzeroMTwoQLIKE}{0.2835}
\newcommand{\mhHOneFourZeroZeroBlkMTwoQLIKE}{0.2731}
\newcommand{\mhHOneFourZeroZeroDMBlkVsAzero}{$-$3.8}
\newcommand{\mhHOneOneThreeZeroAzeroMTwoQLIKE}{0.1632}
\newcommand{\mhHOneOneThreeZeroBlkMTwoQLIKE}{0.1548}
\newcommand{\mhHOneOneThreeZeroDMBlkVsAzero}{$-$3.7}
\newcommand{\mhHOneOneZeroZeroAzeroMTwoQLIKE}{0.1507}
\newcommand{\mhHOneOneZeroZeroBlkMTwoQLIKE}{0.1436}
\newcommand{\mhHOneOneZeroZeroDMBlkVsAzero}{$-$4.2}
\newcommand{\mhHOneThreeThreeZeroAzeroMTwoQLIKE}{0.2410}
\newcommand{\mhHOneThreeThreeZeroBlkMTwoQLIKE}{0.2276}
\newcommand{\mhHOneThreeThreeZeroDMBlkVsAzero}{$-$3.6}
\newcommand{\mhHOneThreeZeroZeroAzeroMTwoQLIKE}{0.2182}
\newcommand{\mhHOneThreeZeroZeroBlkMTwoQLIKE}{0.2095}
\newcommand{\mhHOneThreeZeroZeroDMBlkVsAzero}{$-$2.8}
\newcommand{\mhHOneTwoThreeZeroAzeroMTwoQLIKE}{0.2005}
\newcommand{\mhHOneTwoThreeZeroBlkMTwoQLIKE}{0.1879}
\newcommand{\mhHOneTwoThreeZeroDMBlkVsAzero}{$-$3.9}
\newcommand{\mhHOneTwoZeroZeroAzeroMTwoQLIKE}{0.1830}
\newcommand{\mhHOneTwoZeroZeroBlkMTwoQLIKE}{0.1732}
\newcommand{\mhHOneTwoZeroZeroDMBlkVsAzero}{$-$4.4}
\newcommand{\mhHOneZeroThreeZeroAzeroMTwoQLIKE}{0.1408}
\newcommand{\mhHOneZeroThreeZeroBlkMTwoQLIKE}{0.1338}
\newcommand{\mhHOneZeroThreeZeroDMBlkVsAzero}{$-$4.0}
\newcommand{\mhHOneZeroZeroZeroAzeroMTwoQLIKE}{0.1330}
\newcommand{\mhHOneZeroZeroZeroBlkMTwoQLIKE}{0.1257}
\newcommand{\mhHOneZeroZeroZeroDMBlkVsAzero}{$-$4.9}
\newcommand{\mhHoursBlkBetter}{11}
\newcommand{\mhHoursBlkSig}{11}
\newcommand{\mhHoursTotal}{11}
\newcommand{\mhMOneDMBlkVsAzero}{$-$4.2}
\newcommand{\mhMThreeDMBlkVsAzero}{$-$5.9}
\newcommand{\mhMTwoDMBlkVsAzero}{$-$5.4}
\newcommand{\mhMTwopooledDMBlkVsAzero}{$-$5.3}
\newcommand{\mhNDays}{5{,}587}
\newcommand{\mhNaiveQLIKE}{0.8409}
\newcommand{\mhPooledVsPerHourAzero}{$-$0.3}
\newcommand{\mhPooledVsPerHourBlk}{+0.3}
\newcommand{\mhPostAzeroMOneDMVsNaive}{$-$13.0}
\newcommand{\mhPostAzeroMOneQLIKE}{0.2537}
\newcommand{\mhPostAzeroMThreeDMVsNaive}{$-$13.2}
\newcommand{\mhPostAzeroMThreeQLIKE}{0.2404}
\newcommand{\mhPostAzeroMTwoDMVsNaive}{$-$12.8}
\newcommand{\mhPostAzeroMTwoQLIKE}{0.2511}
\newcommand{\mhPostAzeroMTwopooledDMVsNaive}{$-$12.9}
\newcommand{\mhPostAzeroMTwopooledQLIKE}{0.2506}
\newcommand{\mhPostBlkMOneDMVsNaive}{$-$13.6}
\newcommand{\mhPostBlkMOneQLIKE}{0.2357}
\newcommand{\mhPostBlkMThreeDMVsNaive}{$-$13.7}
\newcommand{\mhPostBlkMThreeQLIKE}{0.2264}
\newcommand{\mhPostBlkMTwoDMVsNaive}{$-$13.4}
\newcommand{\mhPostBlkMTwoQLIKE}{0.2346}
\newcommand{\mhPostBlkMTwopooledDMVsNaive}{$-$13.5}
\newcommand{\mhPostBlkMTwopooledQLIKE}{0.2341}
\newcommand{\mhPostMOneDMBlkVsAzero}{$-$3.2}
\newcommand{\mhPostMThreeDMBlkVsAzero}{$-$4.0}
\newcommand{\mhPostMTwoDMBlkVsAzero}{$-$3.9}
\newcommand{\mhPostMTwopooledDMBlkVsAzero}{$-$3.9}
\newcommand{\mhPostNDays}{2{,}089}
\newcommand{\mhPostNaiveQLIKE}{0.9044}
\newcommand{\mhPostPooledVsPerHourAzero}{$-$1.6}
\newcommand{\mhPostPooledVsPerHourBlk}{$-$1.6}
\fi

\subsection{From One Bar to the Rest of the Session}
\label{sec:multihorizon_methods}

% Methods only: the multi-horizon object, the constructions that reach
% it from a per-bar forecast, the causality convention, the scoring.

The models of this paper issue one forecast per thirty-minute bar. An
intraday user rarely wants only the next bar: at any bar $t$ of a
session the quantity of interest is often the variance still to come
before the close,
\begin{equation}
    RV_{t,T} \;=\; \sum_{t < t_k \le T} RV_{t_k},
    \label{eq:mh_target}
\end{equation}
the sum of the per-bar realized variances from the next bar through the
last regular-hours bar. This is a multi-horizon target whose horizon
shrinks through the day --- eleven bars remain at 10:00, one at 15:00
--- and it must be forecast afresh at every bar, from what is known
then. This subsection fixes how a per-bar forecast is extended to it,
and how the extension is scored. Every construction below is
$\mathcal{F}_t$-measurable: nothing on its right-hand side is formed
after $t$.

\paragraph{The multi-horizon panel.} The target is formed on the paper's
own panel (Section~\ref{sec:data_prep}), restricted to the regular
session, 10:00 to 15:30~ET, twelve bars a day: \mhNDays{} complete
sessions. Entry bars run from 10:00 to 15:00; the 15:30 bar has no
remaining horizon. The per-bar forecasts entering the constructions are
those of the OLS--HAR incumbent (Section~\ref{sec:dp_benchmark}) and the
two-block ridge (Section~\ref{sec:block_tikhonov}), exactly as scored in
Section~\ref{sec:three_block_results}; no forecast is refit for this
purpose.

\paragraph{Constructions.} Let $\hat v_t$ be a model's forecast, issued
at $t$, of the next bar's variance, and let $RV_{\mathrm{open},t}$ be
the variance already realized in the session through bar $t$. Three
nested constructions map the one-step forecast to the remaining
horizon, each an expanding log-linear regression fit separately for
each entry hour $h$ on strictly earlier sessions and applied one session
forward:
\begin{align}
    \text{(M1)}\quad
    \log RV_{t,T} &= a_h + b_h \log \hat v_t + \varepsilon_t,
    \label{eq:mh_m1} \\
    \text{(M2)}\quad
    \log RV_{t,T} &= a_h + b_h \log \hat v_t
                   + d_h \log RV_{\mathrm{open},t} + \varepsilon_t,
    \label{eq:mh_m2} \\
    \text{(M3)}\quad
    \log RV_{t,T} &= a_h + b_h \log \hat v_t
                   + d_h \log RV_{\mathrm{open},t}
                   + e_h \log RV^{(-1)}_{t,T} + \varepsilon_t,
    \label{eq:mh_m3}
\end{align}
where $RV^{(-1)}_{t,T}$ is the previous session's realized variance
over the same remaining horizon. M1 asks how much of the rest of the
session a single next-bar forecast reaches; M2 adds the running
realization of the current session, which is known at $t$ and carries
the session's level so far; M3 adds yesterday's realization of the same
horizon, a seasonal anchor. The regression is what maps thirty minutes
to the rest of the day: its slope on $\log \hat v_t$ is fitted per
hour, because the map from ``the next bar'' to ``the rest of the
session'' differs at 10:00, with the opening bar's turbulence still
ahead, from 14:30, with three bars left. Each fitted value is
back-transformed by the same causal, QLIKE-optimal multiplicative
factor used throughout the paper --- the expanding mean of the
realized-to-forecast ratio over earlier sessions, applied one session
forward --- and no forecast is issued before $63$ sessions of history.
A fourth construction pools the eleven entry hours into one regression
with hour dummies (M2, pooled), so that the gain from hour-specific
slopes can be separated from the gain from the regressors.

\paragraph{Benchmark and scoring.} The horizon-specific benchmark is
the unconditional one: for each entry hour, the expanding mean of
$\log RV_{t,T}$ over earlier sessions, back-transformed the same way
--- the best constant forecast of that horizon with no model in it.
Every construction is scored by QLIKE \citep{Patton2011} on the
remaining variance, one loss per session per entry hour, and reported
per entry hour and pooled across hours (the session mean of the eleven
losses). Two comparisons are made by Diebold--Mariano statistics
\citep{DieboldMariano1995} on paired daily loss differences: each
construction against the horizon benchmark, and, within each
construction, the two-block ridge against the incumbent --- the
multi-horizon counterpart of the one-bar comparison in
Table~\ref{tab:block_ladder}. The sample is also split at
1~January~2016 as a pre-registered era check.
